% !TeX root = ../main.tex
\section{Methodology and Safety}\label{sec:method}
\rubric{Safety and zero-queue baselines}{30}{Safe setup, direct I/O and alignment, and an accurate QD=1
latency table (with \cref{sec:zero-queue}).}
\rubric{Reproducibility and data hygiene}{25}{Configs, raw data, scripts, repeated trials and
environment documentation.}

\subsection{Platform}
\newcommand{\acspinning}{a single I/O-issuing thread, not CPU-pinned (storage I/O is device-bound, not
CPU-bound, below the host-saturation queue depths of \cref{sec:queue-depth})}
\newcommand{\acsenvextra}{%
  SSD       & Micron 2210 \code{MTFDHBA512QFD}, \qty{512}{\giga\byte} QLC NVMe \cite{micron2210brief},
              firmware \code{P6MA001} (SMART log taken with root after the run, \cref{sec:anomalies}) \\
  Interface & PCIe \qty{8.0}{GT/s} $\times$4 (Gen3), current and maximum, from sysfs. Raw ceiling
              $8\,\text{GT/s}\times4\times\tfrac{128}{130}=\qty{3.94}{\giga\byte\per\second}$
              (\qty{3756}{\mebi\byte\per\second}; 128b/130b coding \cite{pcisig_pcie3}), roughly
              \qtyrange{3.2}{3.5}{\giga\byte\per\second} after TLP/DLLP overhead \\
  Block layer & scheduler \code{none}, \code{nr_requests} 1023, 16 I/O queues, volatile write cache
              \enquote{write back} \\
  Target    & disposable files on the SSD's ext4 \code{/} partition (\code{nvme0n1p3}, \code{relatime},
              no online discard) via \code{ACS_SCRATCH}: a \qty{2}{\gibi\byte} test file for the sweeps
              and a \qty{60}{\gibi\byte} file for the original device-state runs; the lean re-run used a
              \qty{4}{\gibi\byte} test file (random offsets over its first \qty{2}{\gibi\byte}) and a
              fill file capped at \qty{40}{\gibi\byte} \\
  Temperature & NVMe composite sensor (\code{hwmon*/temp1_input}), logged per repetition and every
              second in the sustained runs \\
}
%% Measurement environment table (common requirement: measurement discipline).
%% Values come from common/system_info.tex. Before \input-ing this file a project may define
%%   \acspinning   how runs were pinned/isolated
%%   \acsenvextra  extra rows, e.g.  \newcommand{\acsenvextra}{SSD & model, firmware, PCIe link \\}
\providecommand{\acspinning}{\todo{e.g.\ \code{taskset -c 2} on a P-core; isolcpus if used}}
\providecommand{\acsenvextra}{}
\begin{table}[tbp]
  \centering\small
  \caption{Measurement environment.}\label{tab:environment}
  \begin{tabularx}{\linewidth}{@{}>{\raggedright\arraybackslash}p{0.2\linewidth}X@{}}
    \toprule
    CPU               & \sysCPU{} (\sysArch) \\
    SIMD ISA          & \sysISA \\
    Topology          & \sysCores; NUMA: \sysNUMA \\
    SMT               & \sysSMT \\
    Caches            & \sysCaches \\
    Memory            & \sysMemory; DIMMs: \sysDIMM \\
    Pages             & base \sysPageSize; THP: \sysHugePages \\
    OS / kernel       & \sysOS; \sysKernel \\
    Frequency policy  & governor: \sysGovernor; driver: \sysDriver; turbo: \sysTurbo; EPP: \sysEPP \\
    Pinning / isolation & \acspinning; \sysIsolation \\
    Compilers         & \sysGCC; \sysClang \\
    Tools             & \sysPerf{} (perf\_event\_paranoid \sysPerfParanoid); \sysFio; MLC: \sysMLC \\
    GPU (optional)    & \sysGPU; CUDA \sysCUDA \\
    \acsenvextra
    Collected         & \sysCollected \\
    \bottomrule
  \end{tabularx}
\end{table}


The machine also has a SATA HDD (\code{/dev/sda}); each run's \code{env.txt} resolves the scratch
partition to its parent disk and records its \code{rotational} flag, model, firmware, link and
temperature, so a run cannot be silently mislabelled (\cref{sec:ai-use} describes how this check once
failed).

\subsection{Data-loss safeguards}
\begin{itemize}
  \item \textbf{Targets.} Only disposable files in a dedicated scratch directory are written, and both are
        deleted when \code{run.sh} exits, including on interrupt.
  \item \textbf{Gate.} \code{scripts/safety.sh} runs before every experiment and refuses block or character
        devices (also via symlinks), mountpoints, directories, paths outside the scratch directory and
        targets without enough free space; the scratch directory itself may not be \code{/},
        \code{\$HOME}, a filesystem root or foreign-owned. It only \code{stat}s paths
        (\cref{sec:validation}).
  \item \textbf{Benchmark.} \code{blkbench} \code{lstat}s its target, refuses anything but a regular file
        and opens with \code{O_NOFOLLOW}; no raw device is ever opened.
  \item \textbf{Space.} The large device-state file is refused unless \qty{10}{\gibi\byte} stay free.
\end{itemize}

\subsection{I/O engines and configuration}
The primary engine, \code{build/blkbench}, is a small O\_DIRECT \code{libaio} benchmark (the engine FIO
also uses \cite{axboe_fio}). It keeps exactly QD requests in flight in a closed loop and records every
I/O's submit-to-completion latency. FIO 3.42 runs the same points through \code{fio/job.fio}
(\cref{sec:validation}), and \code{ENGINE=fio} routes every sweep through FIO.
\begin{knobs}
  Engine          & \code{libaio}; \code{O_DIRECT}; 4\,KiB-aligned buffers and offsets; 64-bit uniform
                    random offsets (splitmix64) \\
  Timing          & \qty{4}{\second} per repetition (\qty{6}{\second} for tail); write points also stop after
                    \qty{256}{\mebi\byte} (\qty{128}{\mebi\byte} and a \qty{15}{\second} pause in the lean
                    re-run); preconditioning absorbs warm-up \\
  Repetitions     & 5 per point (3 for block-size writes, 4 per engine or span in parts of the lean
                    re-run); median with inter-quartile-range (IQR) bars \\
  Data pattern    & incompressible xorshift64* buffers, fixed per queue slot (we did not test whether the
                    2210 deduplicates) \\
  Latency metric  & per-I/O submit$\to$completion; p50/p95/p99/p99.9 from the full sorted sample \\
  Preconditioning & sequential O\_DIRECT write over the test file before every sequential-read group
                    (\cref{sec:layout}). In the reported \code{all} run the file was otherwise
                    preconditioned only before \code{zero_queue} (the current \code{run.sh} does it before
                    every experiment); the lean re-run did it once, before \code{tail} \\
  Driver          & \code{SSD_Performance_Profiling/scripts/run.sh}; \code{src/blkbench.c} \\
\end{knobs}

\subsection{Execution, run order and idle periods}
The author ran all timed runs in a terminal (governor \code{performance}, turbo and SMT off).
\code{run.sh all} ran in the order of \cref{tab:budget} on 2026-10-21, 10:32--10:54, followed directly by
\code{SUSTAIN_GIB=60 run.sh device_state} until 11:23 (run \code{20261021-103201}); an earlier pilot was
discarded. Nothing in between trimmed or idled the drive, so each experiment inherited the SLC-cache
fill and background work of the previous one, which shaped the write collapses in
\cref{sec:block-size,sec:queue-depth}, the tails in \cref{sec:tail} and \cref{sec:device-state}.

\textbf{Idle period.} An \emph{idle period} means that the drive received no I/O from our scripts for a
stated time (\qtyrange{5}{30}{\minute}), so that the controller could finish background SLC-to-QLC folding
and garbage collection; an \emph{idle drive} is one measured after such a period.

\paragraph{Lean re-run.} The confounded experiments were re-measured that evening with
\code{rerun_clean.sh lean} (run \code{20261021-220614}, 22:06--22:32 EDT; \cref{tab:lean}) after
\code{sudo fstrim -v /} at 21:34:54 and an idle period of at least \qty{30}{\minute}, with the sustained
fill first and stopped soon after the cliff. Lean points replace only the matching configurations of the
original run; captions name each point's run.

\begin{table}[!htbp]
  \centering\small
  \caption{Measured wall-clock time of each experiment in the original 2026-10-21 runs, in execution
  order, and the maximum bytes it can write with the default knobs (byte caps plus
  preconditioning).}\label{tab:budget}
  \begin{tabularx}{\linewidth}{@{}lrrX@{}}
    \toprule
    Experiment & Time & Max written & Notes \\
    \midrule
    \code{zero_queue}   & \qty{1.1}{\minute} & \qty{7}{\gibi\byte}  & includes the initial preconditioning \\
    \code{block_size}   & \qty{9.1}{\minute} & \qty{22}{\gibi\byte} & 48 points; write points 3 reps, \qty{256}{\mebi\byte} cap \\
    \code{rw_mix}       & \qty{0.5}{\minute} & \qty{4}{\gibi\byte}  & \\
    \code{queue_depth}  & \qty{3.7}{\minute} & \qty{11}{\gibi\byte} & read and write, QD 1--256 \\
    \code{tail}         & \qty{2.0}{\minute} & --                   & reads only \\
    \code{pitfall}      & \qty{0.8}{\minute} & \qty{2}{\gibi\byte}  & reads only (plus re-fill) \\
    \code{layout}       & \qty{2.2}{\minute} & \qty{6}{\gibi\byte}  & includes a \qty{4}{\gibi\byte} random overwrite \\
    \code{fio_check}    & \qty{3.3}{\minute} & \qty{8}{\gibi\byte}  & both engines, interleaved \\
    \code{sustained}    & \qty{24.9}{\minute}& \qty{72}{\gibi\byte} & \qty{60}{\gibi\byte} fill $+$ $\le$\qty{8}{\gibi\byte} random $+$ $\le$\qty{4}{\gibi\byte} recovery \\
    \code{working_set}  & \qty{4.0}{\minute} & \qty{5}{\gibi\byte}  & on the \qty{60}{\gibi\byte} file left by \code{sustained} \\
    \bottomrule
  \end{tabularx}
\end{table}

\begin{table}[!htbp]
  \centering\small
  \caption{The lean re-run (\code{rerun_clean.sh lean}, 2026-10-21, 22:06--22:32 EDT) in execution order:
  wall-clock time from the run log, including idle periods, and the most each step could write.}\label{tab:lean}
  \begin{tabularx}{\linewidth}{@{}lrrX@{}}
    \toprule
    Step & Time & Max written & What it re-measured \\
    \midrule
    \code{sustained}   & \qty{1.6}{\minute}  & \qty{6.3}{\gibi\byte} & fill (stopped at \qty{4.3}{\gibi\byte}), \qty{1}{\gibi\byte} random, \qty{60}{\second} idle, \qty{1}{\gibi\byte} recovery \\
    \code{tail}        & \qty{12.1}{\minute} & \qty{4}{\gibi\byte}   & \qty{10}{\minute} idle, fresh \qty{4}{\gibi\byte} precondition, QD 8/32/64/128 reads \\
    \code{working_set} & \qty{1.6}{\minute}  & \qty{0.75}{\gibi\byte} & spans 1, 2, \qty{4}{\gibi\byte}, forward then reverse \\
    \code{block_size}  & \qty{8.5}{\minute}  & \qty{1.1}{\gibi\byte} & QD=32 random reads (8 sizes), \qty{5}{\minute} idle, QD=1 random writes $\geq$\qty{256}{\kibi\byte} \\
    \code{queue_depth} & \qty{1.1}{\minute}  & \qty{2.5}{\gibi\byte} & random writes at QD 4--32 \\
    \code{fio_check}   & \qty{1.0}{\minute}  & \qty{3}{\gibi\byte}   & the three write points, engine order alternating \\
    \bottomrule
  \end{tabularx}
\end{table}
